% Кодировка, язык и воспроизводимые шрифты из TeX Live.
\usepackage{fontspec}
\usepackage{polyglossia}
\setmainlanguage{russian}
\setotherlanguage{english}
\setmainfont{TeX Gyre Pagella}
\setsansfont{TeX Gyre Heros}
\setmonofont{DejaVu Sans Mono}[Scale=MatchLowercase]
% TeX Gyre Pagella/Heros не содержат кириллицу. Polyglossia применяет эти
% воспроизводимые TeX Live fallback-шрифты только для русского текста.
\newfontfamily\cyrillicfont{Tempora}
\newfontfamily\cyrillicfontsf{DejaVu Sans}
\newfontfamily\cyrillicfonttt{DejaVu Sans Mono}[Scale=MatchLowercase]

\usepackage[margin=24mm,headheight=22pt]{geometry}
\usepackage{microtype}
\usepackage{setspace}
\setstretch{1.08}
\usepackage{parskip}
\setlength{\parindent}{0pt}
\setlength{\parskip}{0.55em}
\setlength{\emergencystretch}{3em}

\usepackage{xcolor}
\definecolor{ciaoNavy}{HTML}{17324D}
\definecolor{ciaoBlue}{HTML}{2B78B8}
\definecolor{ciaoPaleBlue}{HTML}{D9ECF8}
\definecolor{ciaoGreen}{HTML}{3D8168}
\definecolor{ciaoAmber}{HTML}{A76813}
\definecolor{ciaoRed}{HTML}{A33A3A}
% Цвет Query совпадает с семантическим цветом draw.io-экспортера CIAO.
\definecolor{ciaoQueryRed}{HTML}{B85450}
\definecolor{ciaoText}{HTML}{202832}
\definecolor{ciaoMuted}{HTML}{5D6874}
\definecolor{ciaoCodeBg}{HTML}{F4F6F8}
\definecolor{ciaoRule}{HTML}{D4DAE0}
\color{ciaoText}

\usepackage{titlesec}
\titleformat{\chapter}[hang]
  {\sffamily\bfseries\color{ciaoNavy}\huge}
  {\makebox[1.7em][l]{\thechapter}}{0pt}{}
\titlespacing*{\chapter}{0pt}{-10pt}{22pt}
\titleformat{\section}
  {\sffamily\bfseries\color{ciaoNavy}\Large}
  {\thesection}{0.65em}{}
\titleformat{\subsection}
  {\sffamily\bfseries\color{ciaoBlue}\large}
  {\thesubsection}{0.55em}{}

\usepackage{fancyhdr}
\pagestyle{fancy}
\fancyhf{}
\fancyhead[L]{\small\sffamily\nouppercase{\leftmark}}
\fancyhead[R]{\small\sffamily\GuideShortTitle}
\fancyfoot[C]{\sffamily\thepage}
\renewcommand{\headrulewidth}{0.4pt}
\renewcommand{\headrule}{\hbox to\headwidth{\color{ciaoRule}\leaders\hrule height \headrulewidth\hfill}}
\fancypagestyle{plain}{%
  \fancyhf{}
  \fancyfoot[C]{\sffamily\thepage}
  \renewcommand{\headrulewidth}{0pt}}

\usepackage{graphicx}
\usepackage{booktabs}
\usepackage{tabularx}
\usepackage{array}
\usepackage{enumitem}
\setlist{itemsep=0.25em,topsep=0.35em}
\usepackage{caption}
\captionsetup{font=small,labelfont={bf,sf},labelsep=period}
\usepackage{tikz}
\usetikzlibrary{arrows.meta,positioning,shapes.geometric,calc}

\usepackage{listings}
\lstdefinestyle{ciaopython}{
  language=Python,
  basicstyle=\ttfamily\small,
  keywordstyle=\color{ciaoBlue}\bfseries,
  commentstyle=\color{ciaoGreen}\itshape,
  stringstyle=\color{ciaoAmber},
  numberstyle=\scriptsize\color{ciaoMuted},
  numbers=left,
  numbersep=9pt,
  frame=single,
  rulecolor=\color{ciaoRule},
  backgroundcolor=\color{ciaoCodeBg},
  xleftmargin=1.3em,
  framexleftmargin=0.8em,
  aboveskip=0.9em,
  belowskip=0.9em,
  showstringspaces=false,
  keepspaces=true,
  columns=fullflexible,
  breaklines=true,
  breakatwhitespace=false,
  postbreak=\mbox{\textcolor{ciaoMuted}{$\hookrightarrow$}\space},
  tabsize=4,
  captionpos=b,
  literate=
    {а}{{а}}1 {б}{{б}}1 {в}{{в}}1 {г}{{г}}1 {д}{{д}}1 {е}{{е}}1
    {ё}{{ё}}1 {ж}{{ж}}1 {з}{{з}}1 {и}{{и}}1 {й}{{й}}1 {к}{{к}}1
    {л}{{л}}1 {м}{{м}}1 {н}{{н}}1 {о}{{о}}1 {п}{{п}}1 {р}{{р}}1
    {с}{{с}}1 {т}{{т}}1 {у}{{у}}1 {ф}{{ф}}1 {х}{{х}}1 {ц}{{ц}}1
    {ч}{{ч}}1 {ш}{{ш}}1 {щ}{{щ}}1 {ъ}{{ъ}}1 {ы}{{ы}}1 {ь}{{ь}}1
    {э}{{э}}1 {ю}{{ю}}1 {я}{{я}}1
    {А}{{А}}1 {Б}{{Б}}1 {В}{{В}}1 {Г}{{Г}}1 {Д}{{Д}}1 {Е}{{Е}}1
    {Ё}{{Ё}}1 {Ж}{{Ж}}1 {З}{{З}}1 {И}{{И}}1 {Й}{{Й}}1 {К}{{К}}1
    {Л}{{Л}}1 {М}{{М}}1 {Н}{{Н}}1 {О}{{О}}1 {П}{{П}}1 {Р}{{Р}}1
    {С}{{С}}1 {Т}{{Т}}1 {У}{{У}}1 {Ф}{{Ф}}1 {Х}{{Х}}1 {Ц}{{Ц}}1
    {Ч}{{Ч}}1 {Ш}{{Ш}}1 {Щ}{{Щ}}1 {Ъ}{{Ъ}}1 {Ы}{{Ы}}1 {Ь}{{Ь}}1
    {Э}{{Э}}1 {Ю}{{Ю}}1 {Я}{{Я}}1
}
\lstset{style=ciaopython}

\usepackage[most]{tcolorbox}
\newtcolorbox{notebox}[1][Примечание]{
  enhanced,breakable,title=#1,colback=ciaoPaleBlue!45,colframe=ciaoBlue,
  fonttitle=\sffamily\bfseries,boxrule=0.7pt,arc=1.2mm}
\newtcolorbox{importantbox}[1][Важно]{
  enhanced,breakable,title=#1,colback=ciaoRed!5,colframe=ciaoRed,
  fonttitle=\sffamily\bfseries,boxrule=0.7pt,arc=1.2mm}
\newtcolorbox{trybox}[1][Попробуйте сами]{
  enhanced,breakable,title=#1,colback=ciaoAmber!7,colframe=ciaoAmber,
  fonttitle=\sffamily\bfseries,boxrule=0.7pt,arc=1.2mm}
\newtcolorbox{resultbox}[1][Результат]{
  enhanced,breakable,title=#1,colback=ciaoGreen!7,colframe=ciaoGreen,
  fonttitle=\sffamily\bfseries,boxrule=0.7pt,arc=1.2mm}

\usepackage{hyperref}
\hypersetup{
  colorlinks=true,
  linkcolor=ciaoBlue,
  urlcolor=ciaoBlue,
  citecolor=ciaoGreen,
  pdftitle={\GuideTitle},
  pdfauthor={\GuideAuthor},
  pdfsubject={\GuideSubtitle},
  pdfcreator={LuaLaTeX}
}
\usepackage[russian]{cleveref}
\crefname{chapter}{глава}{главы}
\Crefname{chapter}{Глава}{Главы}
\crefname{section}{раздел}{разделы}
\Crefname{section}{Раздел}{Разделы}
\crefname{figure}{рисунок}{рисунки}
\Crefname{figure}{Рисунок}{Рисунки}
\crefname{table}{таблица}{таблицы}
\Crefname{table}{Таблица}{Таблицы}
\crefname{lstlisting}{листинг}{листинги}
\Crefname{lstlisting}{Листинг}{Листинги}

\newcommand{\code}[1]{\nolinkurl{#1}}
\newcommand{\RepoPath}[1]{\texttt{\nolinkurl{#1}}}
\newcommand{\RepoLink}[1]{\href{../../#1}{\nolinkurl{#1}}}
\newcommand{\CIAO}{{\sffamily\bfseries CIAO}}
